\documentclass[a4paper,10pt]{article}

\usepackage[top=1.4cm, bottom=1.4cm, left=2cm, right=2cm]{geometry}
\usepackage[T1]{fontenc}
\usepackage{lmodern}
\usepackage{microtype}
\usepackage{parskip}
\usepackage{titlesec}
\usepackage{fontawesome5}
\usepackage{graphicx}
\usepackage[hidelinks, breaklinks=true]{hyperref}

\begin{filecontents*}[overwrite]{refs_UPC.bib}
@article{coppola2019discrete,
  author  = {Coppola, Gennaro and Capuano, Francesco and de Luca, Luigi},
  title   = {Discrete Energy-Conservation Properties in the Numerical Simulation
             of the {Navier--Stokes} Equations},
  journal = {Applied Mechanics Reviews},
  volume  = {71},
  number  = {1},
  pages   = {010803},
  year    = {2019},
  doi     = {10.1115/1.4042820}
}
@article{artiano2025performances,
  author  = {Artiano, Marco and De Michele, Carlo and Capuano, Francesco
             and Coppola, Gennaro},
  title   = {On the Performances of Standard and Kinetic Energy Preserving
             Time-Integration Methods for Incompressible-Flow Simulations},
  journal = {Flow, Turbulence and Combustion},
  volume  = {115},
  pages   = {389--403},
  year    = {2025},
  doi     = {10.1007/s10494-024-00586-8}
}
@misc{neiva2025unfitted,
  author      = {Neiva, Eric and Turlier, Herv{\'e}},
  title       = {Unfitted Finite Element Modelling of Surface-Bulk Viscous Flows
                 in Animal Cells},
  year        = {2025},
  eprint      = {2505.05723},
  eprinttype  = {arxiv},
  eprintclass = {cs.CE},
  doi         = {10.48550/arXiv.2505.05723}
}
@misc{andrews2026stronglyenstrophystableintegratorsincompressible,
      title={Strongly enstrophy-stable integrators for the incompressible Navier-Stokes equations},
      author={Boris D. Andrews and Matin Shams and Patrick E. Farrell},
      year={2026},
      eprint={2609.15520},
      archivePrefix={arXiv},
      url={https://arxiv.org/abs/2609.15520},
      doi={10.48550/arXiv.2609.15520}
}
\end{filecontents*}

\usepackage[backend=biber, style=verbose, maxnames=4, doi=true,
            url=false, isbn=false, eprint=true]{biblatex}
\addbibresource{refs_UPC.bib}
\DeclareFieldFormat{doi}{%
  \mkbibacro{DOI}\addcolon\space
  \textcolor{magenta}{\href{https://doi.org/#1}{\nolinkurl{#1}}}}

\usepackage{amsmath}
\usepackage{amssymb}
\usepackage{xcolor}
\usepackage{textcomp}
\usepackage{array}
\usepackage{colortbl}
\usepackage[most]{tcolorbox}

\pagestyle{empty}
\setlength{\parindent}{1.5em}
\setlength{\parskip}{0pt}

% Body-text scale — matches the CV / Publications documents.
\renewcommand{\normalsize}{\fontsize{12}{13.5}\selectfont}
\renewcommand{\small}{\fontsize{11}{13}\selectfont}

% ── Colours ──────────────────────────────────────────
\definecolor{muted}{gray}{0.45}
\definecolor{linkcol}{HTML}{4C72B0}
\definecolor{midnight}{HTML}{191970}
\definecolor{brick}{HTML}{B5574A}
\definecolor{costBg}{HTML}{F3F5FA}

% ── Section headings — brick bar + small-caps ────────
\newcommand{\sectitlesize}{\fontsize{16}{16}\selectfont}
\newcommand{\sectionbar}{\textcolor{brick}{\rule[-0.05em]{1.1cm}{0.12em}}\hspace{0.45em}}
\titleformat{\section}
  {\color{brick}\scshape\sectitlesize}
  {}{0em}
  {\sectionbar}
\titlespacing{\section}{0pt}{7pt}{5pt}

% ── Subsection headings — brick, plain case ──────────
\newcommand{\subsectitle}[1]{\noindent{\color{brick}\fontsize{15}{17}\selectfont #1}\par\vspace{2pt}}

% ── Linked citation (author, year) ───────────────────────────
\newcommand{\clink}[2]{\textcolor{linkcol}{\href{#1}{#2}}}

% ── Estimated-costs box (dissertation-card style, distinct colour) ──
\newtcolorbox{costbox}{
  enhanced, breakable, colback=costBg, colframe=costBg,
  boxrule=0pt, arc=4pt, left=10pt, right=10pt, top=7pt, bottom=7pt,
  before skip=4pt, after skip=0pt, width=0.85\linewidth, halign=left
}

% ── Cost-table type sizes ─────────────────────────────────────
\newcommand{\costtitlesize}{\fontsize{12}{13}\selectfont}
\newcommand{\costitemsize}{\fontsize{10}{12}\selectfont}
\newcommand{\costtitle}[1]{\textcolor{brick}{\costtitlesize #1}}
\newcommand{\costitem}[1]{\hspace{1em}\costitemsize\textcolor{brick}{$\circ$}\ #1}
\newcommand{\costamt}[1]{\textcolor{midnight}{\costitemsize #1}}
\newcommand{\costamtbold}[1]{\textcolor{midnight}{\textbf{\costtitlesize #1}}}
\newcommand{\budgetnote}[2]{\noindent\footnotesize\textcolor{muted}{\textsuperscript{#1}~#2}\par\vspace{4pt}}

\begin{document}
\normalsize

% ════════════════════════════════════════════════════
%  HEADER
% ════════════════════════════════════════════════════
\noindent
\begin{tabular*}{\textwidth}{@{}l@{\extracolsep{\fill}}r@{}}
{\huge\color{midnight}\scshape Matin Shams}\ \ \textcolor{brick}{\normalsize Motivation Letter}
&
{\Large\textcolor{brick}{\rule[-0.28em]{0.8pt}{1.1em}}\hspace{0.6em}\textcolor{midnight}{%
    \href{mailto:matinshams81@gmail.com}{\faEnvelope[regular]}
    \quad
    \href{https://github.com/matin-shams}{\faGithub}
    \quad
    \href{https://matinshams.me/cv.html?full}{\faGlobe}
  }}
\end{tabular*}

\vspace{20pt}

I am writing to express my interest in the Research Assistant position on lattice Boltzmann methods for wind energy applications at the University of Oldenburg.


% ════════════════════════════════════════════════════
%  ACADEMIC BACKGROUND
% ════════════════════════════════════════════════════
\section*{Education}
I chose to study aerospace engineering. It took me a few years, and a detour through wind farms and black holes, to find my way back. The following subsections trace that path and the research direction it has led to.

\vspace{10pt}
\subsectitle{Undergraduate Studies}
Studying aerospace engineering at TU Delft allowed me to understand the art of simplifying complex
phenomena to solve engineering challenges, yet I found myself drawn to acquire knowledge in a field
of science where I could explore the art of precision and delve further into the \textit{why} question.
My constant curiosity and drive have led me into the cycle of gaining and passing down knowledge
within the scientific community, where I strive to collaborate and truly make an impact. To move
closer to the source of the methods I was using, I sought out research at the interface of
mathematics and engineering. In a project with
\textcolor{midnight}{\href{https://bartripperda.com/}{Bart Ripperda}} and
\textcolor{midnight}{\href{https://smressle.bitbucket.io/index.html}{Sean Ressler}}, through
numerical relativistic fluid simulations of Rayleigh–Taylor instabilities,
we demonstrated that a dense outer accretion disk penetrates an inner dilute magnetosphere, facilitating gas particle movement to the event horizon.
This instilled in me a resolute desire to study advanced mathematics and modern physics, which led me to a minor in astrophysics at the University of British Columbia.

\vspace{5pt}
In the last two years of my undergraduate studies, I worked with \textcolor{midnight}{\href{https://www.tudelft.nl/en/staff/m.i.gerritsma/}{Marc Gerritsma}} on bundle--valued forms in
mimetic spectral element methods, and it was in this project that I first met analysis, PDEs and numerical methods as subjects in their own right rather than as auxiliary tools.
The project introduced me to the idea that a discretisation should not only approximate a PDE, but also retain its structural features at the discrete level.
This was my first serious encounter with the analytic and geometric questions underlying numerical methods.
Motivated by this, I followed the Applied Mathematics bridging programme at TU Delft in my third
year, alongside my regular aerospace coursework and my BSc thesis on regenerative wind farming
(supervised by \textcolor{midnight}{\href{https://www.tudelft.nl/en/ae/organisation/dean-faculty-of-aerospace-engineering/inauguration-henri-werij/carlos}{Carlos Sim\~{a}o Ferreira}} and
\textcolor{midnight}{\href{https://www.tudelft.nl/staff/u.fechner.1/}{Uwe Fechner}}), to prepare for
graduate study in mathematics.

\vspace{10pt}
\subsectitle{Graduate Studies}
This path led me to the MSc in Mathematical Modelling and Scientific Computing at the University of
Oxford. Guided by my academic advisor,
\textcolor{midnight}{\href{https://people.maths.ox.ac.uk/suli/}{Endre S\"uli}}, I completed a few
supererogatory courses on functional analysis and functional analytic methods for PDEs. These
courses granted me the privilege of conducting a dissertation on
\clink{https://matinshams.me/assets/dissertation.pdf}{energy– and
enstrophy–preserving discretisations of the incompressible Navier Stokes equations} under the
supervision of
\textcolor{midnight}{\href{https://pefarrell.org/}{Patrick Farrell}}.
Beyond the dissertation, I implemented physics-informed neural networks in PyTorch with
\textcolor{midnight}{\href{https://www.maths.ox.ac.uk/people/kathryn.gillow}{Kathryn Gillow}}, and
studied high-dimensional L\'evy flights for novelty search with
\textcolor{midnight}{\href{https://www.maths.ox.ac.uk/people/peter.grindrod}{Peter Grindrod}}.

\vspace{10pt}
\subsectitle{2026 (Gap year)}
Alongside the more abstract work, I have wanted at least part of my mathematics to begin at a
table, with children and a piece of chalk. That wish followed me home when I left Oxford and
returned to Iran, where I came across a rural village whose only school survived as a building
and a list of names, with no teacher left to open the door. For a short time I became the missing
teacher, opening the door each morning, writing on a borrowed board, and showing a handful of
children how to line numbers up on a page and make sense of them. Since then, with the help of
local supporters and investors, we have raised enough in funding and commitments to repair the school and guarantee a proper salary for a permanent teacher. With the school now open, my role has become that of someone they can call on.
I can only ask to devote years to the mathematics I love if, somewhere at the
edge of the map, there are children sitting in a classroom with at least the chance to love it too.

\vspace{5pt}
I then co-founded \textit{trew} with
\textcolor{midnight}{\href{https://www.linkedin.com/in/pablo-cembell\%C3\%ADn-mazas-22b68a272/}{Pablo Cembellín}},
whom I met at Oxford, and served as its CTO. The idea was to give firms tamper-evident,
cryptographically verifiable proof of what their AI systems actually decided, a requirement that
regulators are increasingly imposing. Advised by
\textcolor{midnight}{\href{https://www.unibocconi.it/en/faculty/luca-enriques}{Luca Enriques}}, we took the product to a complete MVP. We did not close a funding round, and
are now seeking an acquirer with the help of
\textcolor{midnight}{\href{https://www.eversheds-sutherland.com/en/united-states/people/marzouk-ben}{Ben Marzouk}}.

\vspace{5pt}
I am currently translating
\clink{https://www.cambridge.org/core/elements/abs/medieval-finitism/DF1200CCF8D885032E251462187EEB8D}{Medieval Finitism}
by \textcolor{midnight}{\href{https://www.ms-zarepour.com/}{Mohammad Saleh Zarepour}} into Farsi,
with the help of \textcolor{midnight}{\href{https://www.kaavelajevardi.com/}{Kaveh Lajevardi}}. The
book examines medieval arguments against actual infinity and traces their role
in the later mathematical treatment of infinities of different sizes.

\vspace{5pt}
\section*{Publications \& Preprints}
My dissertation has since been extended into a paper, currently under review at \textit{Foundations of Computational Mathematics}.\footcite{andrews2026stronglyenstrophystableintegratorsincompressible} In this work, we have proposed and analysed a mixed finite element discretisation for the incompressible
Navier--Stokes equations that preserves the
evolution of both energy $\mathcal{K}$ and enstrophy $\mathcal{E}$ at the discrete level.
In two dimensions, the boundedness of enstrophy yields a Reynolds-robust bound on
$\|\nabla\mathbf{u}\|$, providing a natural stabilisation mechanism for under-resolved flows.
Conforming implementations rest on discrete Stokes complexes with
$\mathbf{H}(\mathrm{grad}\,\mathrm{curl})$-- or $H^2$--conforming vorticity spaces, which are available only in limited polynomial degrees and mesh
configurations.
We have therefore introduced two reformulations: a symmetric interior penalty formulation over the
standard discrete de Rham complex, and an
$\mathbf{H}(\mathrm{grad}\,\mathrm{curl})$-- or $H^2$--free reparametrisation via auxiliary variables
using only standard $\mathrm{grad}$--, $\mathrm{curl}$-- and $\mathrm{div}$--conforming spaces. A broad range of boundary conditions can be incorporated.
Numerical experiments confirm the stabilising effect of enstrophy
preservation in both 2D and 3D.

\vspace{10pt}
\subsectitle{Seminars \& Conferences}
Elements of this work have been presented at a number of international conferences and workshops.
I presented it in a minisymposium talk at \clink{https://www.firedrakeproject.org/firedrake_25.html}{Firedrake'25} (University of Leeds) and in a poster at \clink{https://sites.google.com/view/siamukie26/home}{SIAM UKIE '26} (University of Leicester).
\textcolor{midnight}{\href{https://borisandrews.github.io/}{Boris Andrews}} presented it at \clink{https://mfet2025.de/}{ECCOMAS / MFET} (RWTH Aachen University),
\clink{https://cage.ugent.be/acomen2025/}{ACOMEN} (Ghent University), the \clink{https://www.esi.ac.at/events/t3204/}{ESI Programme on Differential Complexes} (University of Vienna),
\clink{https://scicade.org/}{SciCADE} (University of Edinburgh), a workshop on finite element tensor calculus (Tsinghua University),
and the Numerical Analysis Group seminar (University of Oxford).

\vspace{5pt}
\section*{Why This Project}
Wind energy is where my path back to mathematics began: my BSc thesis on regenerative wind farming
was my first sustained encounter with the flows this project targets. My research since has
centred on discretisations of the Navier--Stokes equations that stay stable for under-resolved,
high-Reynolds-number flows, precisely the regime of atmospheric turbulence and turbine wakes.
Lattice Boltzmann methods attack the same problem from the kinetic side, and the question of
which structure a scheme preserves, and how that governs its stability, carries over directly.
Together with my experience of large-scale simulation codes such as \texttt{BHAC} and
\texttt{Athena++}, and of machine-learned turbulence models in \texttt{OpenFOAM}, this project
would let me return to wind energy with the numerical tools I have built since.

\end{document}
